\chapter{全纯函数}
\label{chap:holomorphic-functions}

一元实函数在一点可导，只要求差商沿实轴趋于同一极限。复函数的自变量有两个
实方向，而复导数仍要求差商趋于一个复数，并且极限不得依赖接近方向。这个
局部要求强迫函数满足一组偏微分方程；在后续章节中，它还会导出积分公式、
幂级数展开和很强的唯一性。本章先建立复微分的局部语言。全纯函数一定解析
将在\chapterref{chap:cauchy-integral-formula}由 Cauchy 积分公式推出。

\section{复平面、Riemann 球面和复极限}

把 \(z=x+iy\) 与 \((x,y)\in\R^2\) 对应，得到复平面 \(\C\)。模
\[
 |z|=(x^2+y^2)^{1/2}
\]
给出距离 \(d(z,w)=|z-w|\)，所以开集、紧集、连通性和序列收敛都沿用
\chapterref{chap:euclidean-limits}的定义。非空连通开集称为\term{域}。
若 \(f:E\to\C\)，\(a\) 是 \(E\) 的聚点，则
\[
 \lim_{z\to a}f(z)=L
\]
表示：对每个 \(\varepsilon>0\)，存在 \(\delta>0\)，使 \(z\in E\) 且
\(0<|z-a|<\delta\) 时有 \(|f(z)-L|<\varepsilon\)。写
\(f=u+iv\)，这个极限等价于
\(u\to\operatorname{Re}L\) 与 \(v\to\operatorname{Im}L\)。

加入一点 \(\infty\) 得到扩充复平面
\[
 \widehat{\C}=\C\cup\{\infty\}.
\]
\(\infty\) 的邻域取
\[
 \{\infty\}\cup\{z:|z|>R\}\qquad(R>0).
\]
于是 \(z_n\to\infty\) 当且仅当 \(|z_n|\to\infty\)。映射
\(z\mapsto1/z\) 在约定 \(1/0=\infty\)、\(1/\infty=0\) 后交换 \(0\)
与 \(\infty\)。球面拓扑也由弦距离
\[
 \chi(z,w)=\frac{|z-w|}
 {\sqrt{1+|z|^2}\sqrt{1+|w|^2}},\qquad
 \chi(z,\infty)=\frac1{\sqrt{1+|z|^2}}
\]
给出，其中另约定 \(\chi(\infty,z)=\chi(z,\infty)\)、
\(\chi(\infty,\infty)=0\)。为核实这确为距离，令
\[
 S(x+iy)=\frac{(2x,2y,x^2+y^2-1)}{1+x^2+y^2},
 \qquad S(\infty)=(0,0,1).
\]
分子的长度平方为 \((1+x^2+y^2)^2\)，故像在单位球面上。除北极以外，
球面点 \((X,Y,Z)\) 的逆像为 \((X+iY)/(1-Z)\)，所以 \(S\) 是双射。
直接展开平方得到
\[
 |S(z)-S(w)|^2
 =\frac{4|z-w|^2}{(1+|z|^2)(1+|w|^2)}
 =4\chi(z,w)^2.
\]
一端为无穷点时同样有 \(|S(z)-S(\infty)|=2\chi(z,\infty)\)。
因此弦距离是球面上欧氏距离的一半，特别地满足三角不等式。
有限点处 \(S\) 及上述逆公式连续，给出通常的复平面拓扑；北极附近则有
\[
 B_\chi(\infty,\varepsilon)
 =\{\infty\}\cup
 \{z\in\C:|z|>\sqrt{\varepsilon^{-2}-1}\}
 \qquad(0<\varepsilon<1).
\]
这与约定的无穷点邻域一致。球面是 \(\R^3\) 的闭有界子集，故紧；
\(S\) 把这个紧致度量空间的结构传到 \(\widehat{\C}\)。

\section{复可微与 Cauchy--Riemann 方程}

设 \(D\subseteq\C\) 开，\(f:D\to\C\)，\(z_0\in D\)。若
\[
 f'(z_0)=\lim_{h\to0}\frac{f(z_0+h)-f(z_0)}h
\]
存在，就称 \(f\) 在 \(z_0\)\term{复可微}。这里 \(h\) 从复平面的任意
方向趋于零。若 \(f\) 在 \(D\) 每一点复可微，则称它在 \(D\)\term{全纯}。

写 \(f(x+iy)=u(x,y)+iv(x,y)\)。先令 \(h\in\R\)，再令 \(h=it\)，得到
\[
 f'(z_0)=u_x+iv_x=v_y-iu_y.
\]
因此复可微必有
\begin{equation}
 u_x=v_y,\qquad u_y=-v_x.
 \label{eq:cauchy-riemann-cartesian}
\end{equation}
这就是 Cauchy--Riemann 方程。

两个最简单的例子已经显示实微分与复微分的差别。函数 \(f(z)=z^2\) 满足
\[
 \frac{(z_0+h)^2-z_0^2}{h}=2z_0+h\longrightarrow2z_0,
\]
所以它处处复可微。共轭映射 \(f(z)=\bar z\) 则给出差商
\(\bar h/h\)：沿实方向趋近时它等于 \(1\)，沿纯虚方向趋近时等于
\(-1\)，因此共轭映射在任何一点都不复可微。函数是否复可微取决于增量的
二维行为，而不取决于所考察点的特殊位置。

只在一点验证四个偏导和上述方程并不够。令
\[
 f(x+iy)=
 \begin{cases}
 \displaystyle\frac{x^3}{x^2+y^2}
 +i\frac{y^3}{x^2+y^2},&(x,y)\ne(0,0),\\[0.8em]
 0,&(x,y)=(0,0).
 \end{cases}
\]
在原点有 \(u_x=v_y=1\)、\(u_y=-v_x=0\)，但沿实轴的差商为 \(1\)，
沿 \(y=x\) 的差商为 \(1/2\)。偏导只检测坐标方向，不能替代二维的一阶
线性近似。

\section{全纯函数和实可微结构}

把复函数视为 \(\R^2\to\R^2\) 的映射。若它在 \(z_0\) 实 Fréchet 可微，
则
\[
 f(z_0+h)-f(z_0)=Df(z_0)h+o(|h|).
\]

\begin{theorem}[复可微的实微分判据]
\label{thm:complex-differentiability-real-criterion}
设 \(f=u+iv\) 在 \(z_0\) 实 Fréchet 可微。则 \(f\) 在 \(z_0\) 复可微，
当且仅当 \(u,v\) 满足 Cauchy--Riemann 方程。此时
\[
 Df(z_0)=\begin{pmatrix}a&-b\\ b&a\end{pmatrix},
 \qquad f'(z_0)=a+ib,
\]
其中 \(a=u_x=v_y\)、\(b=v_x=-u_y\)。
\end{theorem}

\begin{proof}
复可微时
\[
 f(z_0+h)-f(z_0)=f'(z_0)h+o(|h|),
\]
所以实微分就是乘以 \(f'(z_0)\) 的映射，其矩阵具有所示形式。反之，实可微
加 Cauchy--Riemann 方程使实微分等于乘以 \(a+ib\)。于是
\[
 \frac{f(z_0+h)-f(z_0)}h
 =a+ib+\frac{o(|h|)}h\longrightarrow a+ib.
\]
\end{proof}

\begin{corollary}[连续偏导判据]
\label{cor:continuous-cr-holomorphic}
若 \(u,v\) 在开集 \(D\) 上具有连续一阶偏导，并在 \(D\) 满足
Cauchy--Riemann 方程，则 \(u+iv\) 在 \(D\) 全纯。
\end{corollary}

这是因为连续偏导保证实 Fréchet 可微。连续性只是一个便于核查的充分条件，
不是复可微定义的一部分。更准确地说，在已知实可微时，Cauchy--Riemann
方程与复可微等价；若只知道四个偏导在一点存在，方程本身并不能补出缺失的
实可微性，上一节的分式函数正是反例。

当 \(f'(z_0)\ne0\) 时，实微分是乘以复数 \(f'(z_0)\) 的映射，因而等于一次
旋转再接一次伸缩；它保持夹角及其定向。若 \(f'(z_0)=0\)，一阶项消失，局部
行为要由更高阶项决定。这一区别以后会分别出现在共形映射和零点重数中。

复导数的线性、乘积、商和链式法则都由差商代数得到。
若 \(f\) 在 \(z_0\) 的整个邻域全纯，并且作为实映射为 \(C^1\)，
且 \(f'(z_0)\ne0\)，则
\[
 \det Df(z_0)=|f'(z_0)|^2>0.
\]
实逆函数定理给出局部 \(C^1\) 逆 \(g\)。缩小邻域后 \(f'\) 处处不为零，
且 \(Dg(w)=[Df(g(w))]^{-1}\) 仍是复线性的，所以 \(g\) 全纯，并有
\[
 g'(w)=\frac1{f'(g(w))}.
\]
只在基点复可微不足以作出这一结论。例如 \(f(z)=z+\bar z^2\) 在零点的
复导数为一，实逆函数定理也给出局部实逆，但它在每个非零点的实微分含有
\(2\bar z\,\bar h\) 项，不是复线性的。若其局部逆全纯，该逆的可逆微分及其逆
都应复线性，便与上述计算矛盾。后面证明全纯函数解析后，邻域 \(C^1\)
条件将自动成立。

\section{复幂级数及收敛圆盘}

对复幂级数
\[
 \sum_{n=0}^{\infty}a_n(z-z_0)^n
\]
令
\[
 R^{-1}=\limsup_{n\to\infty}|a_n|^{1/n},
\]
并采用 \(0^{-1}=\infty\)、\(\infty^{-1}=0\)。根值判别表明
\(|z-z_0|<R\) 时绝对收敛，\(|z-z_0|>R\) 时发散。收敛域是圆盘，而不是
实轴上的区间。根值判别对边界 \(|z-z_0|=R\) 不作结论；同一圆周上的不同
点甚至可能有不同的收敛行为。

\begin{theorem}[复幂级数的逐项求导]
\label{thm:complex-power-series-calculus}
若收敛半径 \(R>0\)，则对每个 \(0<r<R\)，级数在
\(|z-z_0|\le r\) 上绝对一致收敛，并且
\[
 \left(\sum_{n=0}^{\infty}a_n(z-z_0)^n\right)'
 =\sum_{n=1}^{\infty}na_n(z-z_0)^{n-1}
\]
于 \(|z-z_0|<R\) 成立。导数级数具有同一收敛半径。
\end{theorem}

\begin{proof}
取 \(r<\rho<R\)。数项级数 \(\sum|a_n|\rho^n\) 收敛，而
\(|a_n(z-z_0)^n|\le|a_n|\rho^n\) 于闭圆盘成立，所以 M-判别给一致收敛。
导数级数半径不变的证明只涉及系数的绝对值，故实幂级数中的证明原样适用：
对 \(n\ge2\)，
\[
 |na_n|^{1/(n-1)}
 =n^{1/(n-1)}\bigl(|a_n|^{1/n}\bigr)^{n/(n-1)};
\]
其上极限与 \(|a_n|^{1/n}\) 的上极限相同，包括零和无穷的情形。
于是函数级数及导数级数都在较小闭圆盘上一致收敛。

尚须证明所得函数具有复导数，不能直接把实区间的求导定理用于圆盘。
记多项式部分和为 \(S_N\)，导数级数之和为 \(g\)。固定圆盘内的 \(z\)，
取一个较小闭圆盘包含 \(z\) 及充分小的线段 \([z,z+h]\)。
对实参数 \(t\) 应用微积分基本定理，得到
\[
 S_N(z+h)-S_N(z)=h\int_0^1S_N'(z+th)\,dt.
\]
一致收敛允许令 \(N\to\infty\)，故
\[
 \frac{f(z+h)-f(z)}h=\int_0^1g(z+th)\,dt.
\]
函数 \(g\) 连续；当复增量 \(h\to0\) 时右侧趋于 \(g(z)\)，与方向无关。
这就证明逐项求导公式。
\end{proof}

反复求导给
\[
 a_n=\frac{f^{(n)}(z_0)}{n!},
\]
所以已经存在的幂级数表示唯一。任意全纯函数都能这样展开则要等到
Cauchy 积分公式建立之后。

\section{指数、对数和初等全纯函数}

复指数由
\[
 e^z=\sum_{n=0}^{\infty}\frac{z^n}{n!}
\]
定义。Cauchy 乘积给 \(e^{z+w}=e^ze^w\)，逐项求导给
\((e^z)'=e^z\)，并且
\[
 e^{x+iy}=e^x(\cos y+i\sin y).
\]
故 \(e^z\ne0\)，而 \(e^{z+2\pi i}=e^z\)。事实上
\(e^z=e^w\) 当且仅当 \(z-w\in2\pi i\Z\)，所以复指数没有全局单值逆。

在裂平面 \(\C\setminus(-\infty,0]\) 上，每个 \(z\) 唯一写成
\(z=re^{i\theta}\)，其中 \(r>0\)、\(-\pi<\theta<\pi\)。定义
\[
 \operatorname{Log}z=\log r+i\theta.
\]
直接计算实偏导可知它全纯且导数为 \(1/z\)。若域
\(D\subseteq\C\setminus\{0\}\) 上有全纯 \(L\) 满足 \(e^{L(z)}=z\)，
就称 \(L\) 是一个\term{对数分支}。其导数必为 \(1/z\)，两个分支在连通域
上相差 \(2\pi i\) 的整数倍。事实上，若 \(L_1,L_2\) 是两个分支，则
\(e^{L_1-L_2}=1\)；连续函数 \((L_1-L_2)/(2\pi i)\) 取值于离散集
\(\Z\)，在连通域上只能为常数。下一章将用闭曲线积分判定分支是否存在。

正弦和余弦可由幂级数定义，也可写成
\[
 \cos z=\frac{e^{iz}+e^{-iz}}2,\qquad
 \sin z=\frac{e^{iz}-e^{-iz}}{2i}.
\]
它们在复平面上不再有界，例如 \(\cos(iy)=\cosh y\)。选定对数分支后才可
定义 \(z^\alpha=e^{\alpha L(z)}\)；改变分支通常会改变函数。

\section{复可微性的刚性}

\begin{proposition}
\label{prop:zero-derivative-holomorphic-constant}
若域 \(D\) 上的全纯函数满足 \(f'\equiv0\)，则 \(f\) 为常数。
\end{proposition}

\begin{proof}
每一点都有一个闭含于 \(D\) 的小矩形。矩形中任意两点可用水平和竖直线段
连接；沿每条线段应用一元链式法则与“导数为零则常值”，得到端点函数值相同。
因此 \(f\) 局部常值。一个局部常值函数在连通域上只能取一个值。
\end{proof}

若全纯函数只取实值，Cauchy--Riemann 方程给其实部的两个偏导为零，所以它
为常数。像落在某一条仿射直线中的情形经平移、旋转后相同。若 \(|f|\) 恒定且
非零，对 \(u^2+v^2=c^2\) 求偏导并结合 Cauchy--Riemann 方程，也得到
四个偏导全为零。幂级数还显示：若其和在中心的所有阶导数都为零，则所有系数
为零。后面的恒等定理会把这一性质推广到全部全纯函数。

\section{Wirtinger 导数和实微分表示}

对实可微的 \(f=u+iv\)，定义
\[
 f_z=\frac12(f_x-if_y),\qquad
 f_{\bar z}=\frac12(f_x+if_y).
\]
展开后
\[
 f_z=\frac12[(u_x+v_y)+i(v_x-u_y)],\qquad
 f_{\bar z}=\frac12[(u_x-v_y)+i(v_x+u_y)].
\]
所以 Cauchy--Riemann 方程等价于 \(f_{\bar z}=0\)，此时 \(f'=f_z\)。

每个实线性映射 \(L:\C\to\C\) 唯一写成
\[
 L(h)=ah+b\bar h.
\]
应用于实微分得到
\begin{equation}
 Df(z)h=f_z(z)h+f_{\bar z}(z)\bar h.
 \label{eq:wirtinger-real-differential}
\end{equation}
乘以复数的映射恰好对应 \(b=0\)，所以实可微函数复可微当且仅当其实微分
为复线性。例如 \(f(z)=|z|^2=z\bar z\) 满足
\(f_z=\bar z\)、\(f_{\bar z}=z\)，故只在原点复可微，并不在任何域上全纯。

\section*{习题}
\addcontentsline{toc}{section}{习题}
\subsection*{A组}
\begin{enumerate}[label=\textbf{\thechapter.\arabic*.}]
\exerciseitem{complex-spherical-convergence}设 \(z_n\in\widehat{\C}\)，沿用 \(1/0=\infty\)、\(1/\infty=0\)。证明 \(z_n\to\infty\) 当且仅当 \(1/z_n\to0\)，并验证弦距离给出的 \(\infty\) 邻域与正文定义相同。
\exerciseitem{complex-limit-components}从定义证明复函数极限存在当且仅当实部、虚部极限分别存在，并推出复乘法保持极限。
\exerciseitem{cr-point-counterexample}核对正文中分式函数在原点的四个偏导、Cauchy--Riemann 方程和两条路径上的差商。
\exerciseitem{complex-linear-real-matrix}证明实线性映射 \(L:\C\to\C\) 为复线性，当且仅当其矩阵为 \(\begin{psmallmatrix}a&-b\\b&a\end{psmallmatrix}\)。
\end{enumerate}
\subsection*{B组}
\begin{enumerate}[label=\textbf{\thechapter.\arabic*.},resume]
\exerciseitem{continuous-cr-sufficiency}从实 Fréchet 可微定义完整证明“连续偏导加 Cauchy--Riemann 方程”推出全纯。
\exerciseitem{holomorphic-real-line-image}证明域上的全纯函数若像包含在某一条仿射直线中，则函数为常数。
\exerciseitem{holomorphic-constant-modulus}证明域上的全纯函数若模恒定，则函数为常数，并分别处理常数模为零和非零的情形。
\exerciseitem{complex-power-series-derivatives}设 \(f(z)=\sum_{n\ge0}a_n(z-z_0)^n\) 的半径为 \(R>0\)。证明 \(f^{(k)}(z_0)=k!a_k\)，并说明级数表示的唯一性。
\exerciseitem{complex-exponential-fibres}证明 \(e^z=e^w\) 当且仅当 \(z-w\in2\pi i\Z\)，并求复指数在一点的全部逆像。
\exerciseitem{principal-logarithm-derivative}从 \(u(x,y)=\frac12\log(x^2+y^2)\) 和主辐角的偏导出发，证明主值对数全纯且导数为 \(1/z\)。
\end{enumerate}
\subsection*{C组}
\begin{enumerate}[label=\textbf{\thechapter.\arabic*.},resume]
\exerciseitem{wirtinger-real-linear-decomposition}证明每个实线性 \(L:\C\to\C\) 唯一写成 \(L(h)=ah+b\bar h\)，并推导式\eqref{eq:wirtinger-real-differential}。
\exerciseitem{polar-cauchy-riemann}设 \(f=u+iv\) 在不含原点的开集上 \(C^1\)。
在局部极坐标片 \(r>0\) 中证明：\(f\) 全纯当且仅当
\(u_r=r^{-1}v_\theta\)、\(v_r=-r^{-1}u_\theta\)。用于检验 \(z\mapsto z^n\)，其中 \(n\in\Z\)。
\optionalexerciseitem{local-holomorphic-inverse}设 \(f\) 在 \(z_0\) 邻域 \(C^1\) 且全纯，\(f'(z_0)\ne0\)。由实逆函数定理证明局部逆全纯，并求其导数。
\end{enumerate}
